\documentclass[10pt,a4paper]{amsart}
\usepackage[margin=19mm]{geometry}
\usepackage{amsmath,amssymb,mathtools}
\usepackage[scr=rsfs,cal=euler]{mathalfa}
\usepackage{xfrac}
\usepackage[unicode,hidelinks,bookmarks=false]{hyperref}
\newcommand{\ie}{i.e.\ }
\newcommand{\cf}{cf.\ }
\newcommand{\resp}{respectively\ }
\newcommand{\Subsection}{\subsection}
\providecommand{\nobreakdash}{\nobreak\hbox{-}}
\setlength{\emergencystretch}{2em}
\multlinegap0pt
\usepackage{bm}
\usepackage{bbm}
\usepackage{dsfont}
\usepackage{xspace}
\usepackage{cancel}
\usepackage{mathtools}
\usepackage{amsthm}
\makeatletter
\@ifundefined{theorem}{\newtheorem{theorem}{Theorem}}{}
\@ifundefined{proposition}{\newtheorem{proposition}[theorem]{Proposition}}{}
\makeatother
\newcommand{\Dfun}{\mathscr D}
\newcommand{\LD}{\mathcal L_D}
\newcommand{\R}{\mathbb R}
\newcommand{\abs}[1]{\left\lvert #1\right\rvert}
\newcommand{\ip}[2]{\langle #1,#2\rangle}
\newcommand{\norm}[1]{\left\lVert #1\right\rVert}
\newcommand{\rhoE}{\rho}
\newcommand{\tr}{\operatorname{tr}}
\title[Gromov's Euclidean Endpoint $C^0$ Rigidity for the Positive ]{Gromov's Euclidean Endpoint $C^0$ Rigidity for the Positive Mass Theorem}
\author{Jiangcheng You, Heng Zhang}
\date{Source: arXiv:2606.15372v2 (CC BY 4.0)}
\begin{document}
\maketitle

\noindent\emph{Source and licence.} Gromov's Euclidean Endpoint $C^0$ Rigidity for the Positive Mass Theorem. Version arXiv:2606.15372v2 (CC BY 4.0), distributed under the Creative Commons Attribution 4.0 International licence, as recorded in the arXiv licence metadata for this exact version and shown on the abstract page https://arxiv.org/abs/2606.15372v2. Full rights notice: licenses/arxiv-2606.15372-CC-BY-4.0.txt.\par\smallskip

\noindent\textit{Editorial scope.} Counted unit: the Proposition whose source label is \texttt{prop:D-quadratic} in the article (source0.tex lines 959--985, source label \texttt{prop:D-quadratic}), delivered together with the article's own complete original proof of it (source0.tex lines 987--1099, one \texttt{proof} environment of 113 lines). No mathematics is written, repaired or re-derived: statement and proof are the article's own text line for line. Required context retained uncounted: none. Every definition, display and label of those lines is retained unchanged. Omitted from this package and not redistributed: 1--958, 986--986, 1100--1727. Nothing inside a retained excerpt is omitted or altered: every retained body is delivered exactly as the article's lines are written, so no omission receipt of any kind accompanies this package. Citations: the retained excerpts carry no citation command at all, so no bibliography entry of the article is transcribed and no reference is redistributed. Reference adaptation: none was needed and none was made; every reference inside the delivered unit resolves inside it, so no unresolved reference or citation is produced. Printed slips kept literal: no printed word, symbol, hypothesis or number of the counted statement or of its proof was altered. Layout and class: the article's own class is \texttt{amsart} with packages including geometry, amsmath, amssymb, amsfonts, amsthm, mathtools, mathrsfs, enumitem, hyperref; the wrapper is the lane's pinned \texttt{amsart} editorial wrapper at 10pt on A4, which restates the theorem-like environments the retained text needs and adds the article's own macro definitions that the retained text needs (\texttt{\textbackslash Dfun}, \texttt{\textbackslash LD}, \texttt{\textbackslash R}, \texttt{\textbackslash abs}, \texttt{\textbackslash ip}, \texttt{\textbackslash norm}, \texttt{\textbackslash rhoE}, \texttt{\textbackslash tr}), with extra wrapper packages (bm, bbm, dsfont, xspace, cancel, mathtools). The delivered numbering is the unit's own. Assets: no retained span contains a figure environment, a graphics-inclusion command, a PDF-page inclusion, a table or a TikZ picture; no image file is copied into this package, the package declares no asset dependency and no image file is redistributed with it. Licence: version arXiv:2606.15372v2 of this article is distributed under the Creative Commons Attribution 4.0 International licence, as recorded in the arXiv licence metadata for this exact version and shown on the abstract page https://arxiv.org/abs/2606.15372v2. A full rights notice is delivered at licenses/arxiv-2606.15372-CC-BY-4.0.txt. Characters: every retained excerpt is pure ASCII, so the delivered engine, XeLaTeX with its default Latin Modern text font, reports no missing character.
\begin{proposition}\label{prop:D-quadratic}
Fix $p>3$.  There exist $\delta_0>0$ and $C<\infty$, depending on $\Omega$, $p$, and $\psi$, such that if $k$ is a continuous symmetric two-tensor on $\Omega$ and $v\in W^{1,p}(\Omega)$ satisfy
$$
        \norm{k}_{C^0(\Omega)}+
        \norm{v}_{W^{1,p}(\Omega)}\le\delta_0,
$$
then
\begin{equation}\label{eq:D-quadratic}
        \Dfun(\delta+k,\rhoE+v)
        =\Dfun(\delta,\rhoE)+\LD(k,v)+\mathcal R(k,v),
\end{equation}
where
\begin{align}\label{eq:L-def}
        \LD(k,v)=
        &\int_\Omega\phi'(\rhoE)\abs{\nabla\rhoE}^3v\,dy
        +3\int_\Omega\phi(\rhoE)\abs{\nabla\rhoE}\ip{\nabla\rhoE}{\nabla v}\,dy\notag\\
        &+\int_\Omega\phi(\rhoE)
        \left[\frac12(\tr k)\abs{\nabla\rhoE}^3
        -\frac32\abs{\nabla\rhoE}\,k(\nabla\rhoE,\nabla\rhoE)\right]dy,
\end{align}
and
\begin{equation}\label{eq:D-rem-bound}
        \abs{\mathcal R(k,v)}
        \le C\left(\norm{k}_{C^0(\Omega)}+
        \norm{v}_{W^{1,p}(\Omega)}\right)^2.
\end{equation}
\end{proposition}

\begin{proof}
Set
$$
        \alpha=\norm{k}_{C^0(\Omega)},\qquad
        \beta=\norm{v}_{W^{1,p}(\Omega)},\qquad
        \theta=\alpha+\beta.
$$
We may assume $\theta\le1$ by decreasing $\delta_0$.  Since $p>3$,
$W^{1,p}(\Omega)\hookrightarrow C^0(\Omega)$, so
\begin{equation}\label{eq:Sobolev-small-v4}
        \norm{v}_{C^0}+\norm{\nabla v}_{L^r}
        \le C_r\beta,
        \qquad 1\le r\le p.
\end{equation}
The annulus $\Omega$ is separated from the origin, hence $\rhoE$ and $\nabla\rhoE$ are smooth and bounded on $\Omega$.  Decreasing $\delta_0$ if necessary, $\delta+k$ is uniformly positive definite and $\rhoE+v$ stays in a compact subset of $(0,\infty)$ on which all derivatives of $\phi$ are bounded.

First keep the metric Euclidean.  Taylor's theorem gives the pointwise estimate
\begin{equation}\label{eq:phi-remainder-v4}
        \phi(\rhoE+v)=\phi(\rhoE)+\phi'(\rhoE)v+R_\phi,
        \qquad |R_\phi|\le C\norm{v}_{C^0}^2.
\end{equation}
The elementary cubic inequality
\begin{equation}\label{eq:cubic-estimate-v4}
        \left||X+Y|^3-|X|^3-3|X|\ip{X}{Y}\right|
        \le C(|X||Y|^2+|Y|^3)
\end{equation}
with $X=\nabla\rhoE$ and $Y=\nabla v$ gives
\begin{equation}\label{eq:cubic-rho-v4}
        |\nabla(\rhoE+v)|^3
        =|\nabla\rhoE|^3
        +3|\nabla\rhoE|\ip{\nabla\rhoE}{\nabla v}
        +R_\nabla,
        \qquad
        |R_\nabla|\le C(|\nabla v|^2+|\nabla v|^3).
\end{equation}
Multiplying \eqref{eq:phi-remainder-v4} and \eqref{eq:cubic-rho-v4}, the terms linear in $v$ are exactly
$$
        \int_\Omega \phi'(\rhoE)|\nabla\rhoE|^3v\,dy
        +3\int_\Omega \phi(\rhoE)|\nabla\rhoE|
        \ip{\nabla\rhoE}{\nabla v}\,dy.
$$
All remaining function-only terms are quadratic.  Indeed, using \eqref{eq:Sobolev-small-v4},
$$
        \int_\Omega |R_\phi|\,dy\le C\beta^2,
        \qquad
        \int_\Omega |R_\nabla|\,dy
        \le C(\norm{\nabla v}_{L^2}^2+\norm{\nabla v}_{L^3}^3)
        \le C\beta^2,
$$
and the mixed first-order product obeys
$$
        \int_\Omega |v|\,|\nabla v|\,dy
        \le \norm{v}_{C^0}\norm{\nabla v}_{L^1}
        \le C\beta^2.
$$
The remaining products are also quadratic:
$$
        \int_\Omega |R_\phi|(1+|\nabla v|+|\nabla v|^2+|\nabla v|^3)\,dy
        +\int_\Omega |v|\,|R_\nabla|\,dy
        \le C\beta^2.
$$
Thus
\begin{equation}\label{eq:function-only-quadratic-v4}
        \left|\Dfun(\delta,\rhoE+v)-\Dfun(\delta,\rhoE)-\LD(0,v)\right|
        \le C\beta^2.
\end{equation}

Now include the metric perturbation.  For a positive definite matrix $H$ near $I$ and $\xi\in\R^3$, set
$$
        \Psi(H,\xi)=\ip{H^{-1}\xi}{\xi}^{3/2}\sqrt{\det H}.
$$
Uniform Taylor expansion at $I$ gives
\begin{equation}\label{eq:Psi-Taylor-v4}
        \Psi(I+S,\xi)=|\xi|^3+\mathcal M(S,\xi)+R_{\mathcal M}(S,\xi),
\end{equation}
where
$$
        \mathcal M(S,\xi)=\frac12(\tr S)|\xi|^3
        -\frac32|\xi|S(\xi,\xi),
        \qquad
        |R_{\mathcal M}(S,\xi)|\le C|S|^2|\xi|^3.
$$
Applying this with $S=k$ and $\xi=\nabla(\rhoE+v)$, the pure metric remainder satisfies
$$
        \int_\Omega |\phi(\rhoE+v)R_{\mathcal M}(k,\nabla(\rhoE+v))|\,dy
        \le C\alpha^2\int_\Omega (1+|\nabla v|^3)\,dy
        \le C\alpha^2.
$$
It remains to replace the linear metric term
$\phi(\rhoE+v)\mathcal M(k,\nabla(\rhoE+v))$ by
$\phi(\rhoE)\mathcal M(k,\nabla\rhoE)$.  Since $\nabla\rhoE$ is bounded on $\Omega$, and since $\mathcal M$ is linear in $k$ and cubic in $\xi$,
\begin{equation}\label{eq:M-Lipschitz-v4}
        |\mathcal M(k,\nabla(\rhoE+v))-\mathcal M(k,\nabla\rhoE)|
        \le C\alpha(|\nabla v|+|\nabla v|^2+|\nabla v|^3).
\end{equation}
Moreover,
$$
        |\phi(\rhoE+v)-\phi(\rhoE)|\le C|v|.
$$
Consequently
$$
\begin{aligned}
        &\int_\Omega \left|
        \phi(\rhoE+v)\mathcal M(k,\nabla(\rhoE+v))
        -\phi(\rhoE)\mathcal M(k,\nabla\rhoE)
        \right|dy                                                \\
        &\qquad \le C\alpha\int_\Omega
        (|v|+|\nabla v|+|\nabla v|^2+|\nabla v|^3)dy
        \le C\alpha\beta\le C\theta^2.
\end{aligned}
$$
Combining this estimate with \eqref{eq:function-only-quadratic-v4} and the pure metric remainder, the first variation is precisely the expression in \eqref{eq:L-def}, and every leftover term is bounded by $C\theta^2$.  This proves \eqref{eq:D-rem-bound}.
\end{proof}

\end{document}
